

% \section{Motivation}
% A large body of prior work evaluates specification quality by directly verifying whether a concrete implementation satisfies a given specification using tools such as Dafny or Frama-C. These approaches are effective for checking targeted correctness properties with respect to specific verification goals. However, they are inherently goal-dependent: a specification that suffices to discharge a particular proof obligation may still be overly permissive. As a result, such specifications often over-approximate the intended semantics and fail to faithfully characterize the full behavior of the program.

% To move beyond goal-specific verification, \textsc{Clever} proposes a more stringent evaluation protocol in which an LLM is required to generate both a specification and a corresponding implementation, and then prove that the specification accepts the implementation while being isomorphic to a human-written ground-truth specification. This design aims to assess not only whether a specification is verifiable, but also whether it is semantically aligned with the intended program behavior. However, \textsc{Clever} does not directly address the over-approximation problem: requiring isomorphism to a ground-truth specification only ensures structural equivalence, but does not guarantee that the specification tightly captures the intended semantics or rules out unintended behaviors. Moreover, the ground-truth specifications in \textsc{Clever} are themselves not formally verified against any reference implementation, meaning their correctness relies solely on human judgment. If the ground-truth specification is overly permissive or even erroneous, the isomorphism check will simply propagate these weaknesses.

% In practice, however, we find the \textsc{Clever} setting to be prohibitively difficult. Even for moderately complex programs, requiring LLMs to simultaneously construct a specification that subsumes its generated implementation and establish or verify an explicit isomorphism between two independently written specifications far exceeds the current capabilities of both LLMs and automated proof assistants. This tight coupling of specification synthesis, implementation generation, and cross-specification equivalence checking introduces substantial verification overhead and limits scalability. Consequently, the excessive difficulty of this setting diminishes its discriminative power; when the verification bottleneck becomes insurmountable for nearly all cases, CLEVER fails to effectively differentiate the quality of various generated specifications.

% Moreover, \textsc{Clever} relies on Lean as both the specification and implementation language. While Lean is well suited for expressing and proving mathematical properties, it is not designed for writing realistic program implementations, nor is it used as an implementation language in industrial software development. More fundamentally, the rigid separation between specification and implementation within a single ITP language is itself problematic. In mathematical terms, specifications and implementations are closely related—both describe input-output relations, differing primarily in their intended use rather than their formal structure. Specifications may naturally contain executable components (e.g., reference implementations or computational witnesses), and enforcing a strict syntactic boundary between the two is neither necessary nor always meaningful. Furthermore, using an ITP language to write implementations conflates the role of verification with that of software engineering; realistic implementations should be written in mainstream programming languages that reflect actual development practices.

% % Motivated by these challenges, we adopt Rocq as the underlying verification framework. As the de facto standard for software verification (e.g., CompCert), Rocq provides a more suitable foundation for evaluating the semantic correctness and strength of generated specifications, while allowing specifications to be assessed based on how precisely they capture program behavior rather than on rigid syntactic separation.

% In light of these limitations, we propose a different evaluation framework. We adopt Rocq as the underlying verification framework due to its status as the de facto standard for software verification (e.g., CompCert). Unlike \textsc{Clever}, our approach does not rely on implementations written in any ITP language; instead, we derive comprehensive test cases from reference implementations written in mainstream programming languages. To mitigate the over-approximation problem, we assess whether a specification is tight enough to reject incorrect implementations, rather than merely accepting correct ones.

% Ensuring the correctness of software systems is a fundamental challenge in computer science, particularly for safety-critical applications. Interactive Theorem Provers (ITPs) such as Rocq~\cite{huet1997coq}, Lean~\cite{de2015lean}, and Isabelle~\cite{paulson1994isabelle} provide the most rigorous means of achieving this goal, enabling developers to formally specify system behavior and construct machine-checkable proofs that verify implementation compliance. However, this high level of assurance comes at a significant cost: writing formal specifications and conducting rigorous proofs require deep expertise and are extraordinarily labor-intensive, often spanning years for complex systems such as the seL4 verified microkernel~\cite{klein2009sel4}.

% The emergence of Large Language Models (LLMs) has sparked considerable interest in auto-formalization and neural theorem proving. While the mathematical community has achieved promising results integrating LLM-based reasoning into ITPs, demonstrated through notable performance on established benchmarks (e.g., miniF2F) and advances in automated proof synthesis, these advances do not directly address the distinct challenges of program verification, which involves flexible data types and control flow constructs absent in pure mathematical reasoning.

% Existing approaches for evaluating specification quality typically verify whether an implementation satisfies a given specification using tools such as Dafny or Frama-C. However, these approaches are inherently goal-dependent: a specification that suffices to discharge a particular proof obligation may still be overly permissive, failing to faithfully characterize the full behavior of the program. CLEVER~\cite{clever} introduced the first curated benchmark for formally verified code generation based on HumanEval~\cite{chen2021humaneval}, proposing a more stringent evaluation protocol requiring LLMs to generate both specifications and implementations while proving semantic equivalence to ground-truth specifications. However, CLEVER faces several critical limitations. First, it does not directly address the over-approximation problem---requiring isomorphism to a ground-truth specification ensures structural equivalence but does not guarantee that specifications tightly capture intended semantics or rule out unintended behaviors. Second, the ground-truth specifications themselves are not formally verified against any reference implementation, meaning their correctness relies solely on human judgment. Third, the setting proves prohibitively difficult in practice: simultaneously constructing specifications, implementations, and cross-specification equivalence proofs far exceeds current capabilities of both LLMs and automated proof assistants. Even advanced models achieve only marginal success rates (1.242\% and 5.559\% on subtasks, 0.621\% end-to-end). When the verification bottleneck becomes insurmountable for nearly all cases, the benchmark fails to effectively differentiate the quality of various generated specifications. Finally, CLEVER's reliance on Lean as both specification and implementation language conflates verification with software engineering. While Lean is well suited for expressing and proving mathematical properties, it is not designed for writing realistic program implementations. Realistic implementations should be written in mainstream programming languages that reflect actual development practices.

% In light of these limitations, we propose a novel evaluation framework focused on assessing the quality of LLM-generated formal specifications using Rocq---the de facto standard for software verification (e.g., CompCert). Unlike CLEVER, our approach does not rely on implementations written in any ITP language; instead, we derive comprehensive test cases from reference implementations written in mainstream programming languages. We adopt a practical, test case-based verification methodology: rather than requiring full formal proofs of equivalence, we verify that a set of positive (input, output) test cases satisfies the specifications. To strengthen the assessment and mitigate the over-approximation problem, we also employ negative test cases to check that incorrect behaviors are properly ruled out---assessing whether a specification is tight enough to reject incorrect implementations, rather than merely accepting correct ones.

% The contributions of this paper are as follows:
% \begin{itemize}
%     \item We construct a curated dataset of manually written Rocq specifications for all 164 HumanEval problems, serving as a human baseline for specification quality assessment. 
%     \item We propose a novel evaluation framework with a suite of fine-grained metrics for assessing the quality of LLM-generated formal specifications. We conduct comprehensive experiments across a wide range of state-of-the-art LLMs and compare their performance against the human baseline.
%     \item We perform an in-depth analysis of the evaluation results, addressing several key research questions regarding LLMs' capabilities in formal specification generation.
% \end{itemize}


Formal verification provides the strongest guarantees of software correctness by enabling machine-checkable proofs that an implementation satisfies its intended behavior. Interactive theorem provers (ITPs) such as Rocq~\cite{huet1997coq}, Lean~\cite{de2015lean}, and Isabelle~\cite{paulson1994isabelle} have been successfully applied to the verification of large-scale, safety-critical systems, including certified compilers and operating system kernels. Despite these successes, the practical adoption of formal verification remains fundamentally constrained by the cost of writing formal specifications. Constructing specifications that are both semantically precise and sufficiently restrictive to rule out unintended behaviors requires expert knowledge and substantial manual effort, often exceeding the effort required to verify the implementation itself.


Recent advances in large language models (LLMs) have demonstrated that neural models can meaningfully assist formal reasoning tasks. In particular, LLMs have achieved notable success in formal mathematical theorem proving~\cite{xin2024deepseek,wang2024lego,xin2025deepseek,ren2025deepseek,chen2025seed,lin2025goedel,wang2025kimina,hubert2025olympiad}, where models generate or complete proofs that are mechanically checked by proof assistants. These results indicate that LLMs can internalize non-trivial logical structure and interact effectively with formal languages. Building on this progress, recent work has begun to transfer these capabilities to program verification, where LLMs are used for tasks such as invariant inference~\cite{kamath2023finding,chakraborty2023ranking,pirzada2024llm,cao2025clause2inv}, translating natural language into formal specifications~\cite{cosler2023nl2spec,endres2024can,cao2025informal,fang2025enhancing}, and program synthesis~\cite{first2023baldur,misu2024towards}. This progression naturally raises a further question: \emph{how powerful are LLMs at generating formal program specifications?} Unlike theorem proving or implementation synthesis, specification generation requires translating informal intent---expressed through natural language descriptions or reference implementations---into precise logical predicates that characterize program behavior.


% A central challenge in answering this question lies in how specification quality is evaluated. 

A central challenge in answering this question lies in how to evaluate specification quality at scale. Relying solely on human experts is not only inefficient but also cannot guarantee correctness, as even experts may overlook subtle semantic errors. Most existing work in program verification evaluates specifications indirectly by checking whether a given implementation satisfies them. In these works, specifications take the form of verification objectives, or error-checking conditions, and are designed primarily to discharge a specific verification goal rather than to fully characterize program behavior. While effective for validating implementations against targeted properties or ruling out particular classes of errors, this paradigm provides only a partial view of specification quality: a specification may be sufficient to verify a particular implementation or objective yet remain overly permissive, failing to rule out unintended behaviors or alternative incorrect implementations. In other words, traditional verification pipelines are fundamentally oriented toward establishing $\textit{implementation} \models \textit{specification}$, but offer limited ability to assess whether a specification itself precisely captures the intended semantics of the program.

Recent benchmarks have attempted to strengthen specification evaluation by imposing stricter correctness criteria. \textsc{Clever}~\cite{thakur2025clever}, for example, proposes an end-to-end benchmark that requires models to generate specifications together with implementations and proofs of equivalence against a claimed ground-truth specification, all expressed within a theorem-proving language. While conceptually appealing, this design introduces substantial practical challenges. 

First, while equivalence proofs can in principle avoid over-approximation, they reintroduce dependence on human-curated ground-truth specifications—resources that demand substantial effort to produce and may themselves introduce bias or inconsistency. Second, even verifying that a generated implementation satisfies a generated specification in a fully formal setting is already highly challenging. Requiring, in addition, full equivalence proofs between non-trivial specifications further increases the proof burden, resulting in very low success rates even for state-of-the-art models. Third, because specification generation, implementation generation, and formal verification are tightly coupled in a single pipeline, failures are difficult to diagnose: it is often unclear whether an unsuccessful result reflects an incorrect specification, an incorrect implementation, or limitations of the proving process itself. As a result, although \textsc{Clever} enforces a strong notion of correctness, its evaluation signal is sparse and hard to interpret, limiting its usefulness for assessing specification quality in isolation.



These limitations suggest the need for an evaluation methodology that can assess specification quality within a tractable and interpretable formal reasoning framework. In this work, we argue that specification evaluation should focus on what can be reliably established through formal proofs rather than relying solely on all-or-nothing semantic equivalence checking.

Motivated by these considerations, we propose \textsc{Coins} (\textbf{CO}q-based \textbf{IN}stantiated \textbf{S}pecification evaluation), an evaluation framework for assessing LLM-generated formal specifications using the Rocq proof assistant. The name reflects the core idea of our approach: specifications are evaluated by instantiating them on concrete behaviors and reasoning about the resulting proof obligations in Rocq. We adopt Rocq because it is one of the most widely used proof assistants for program verification~\cite{leroy2016compcert,cao2018vst,zhou2024vst,wu2025qcp}, with a long history of successful applications in large-scale verified systems such as CompCert~\cite{leroy2016compcert}. We chose Rocq over Lean because Rocq has a longer and more established track record in program verification, making it a more natural fit for evaluating program specifications.

Using \textsc{Coins}, we conduct a large-scale empirical study of formal specification generation in Rocq on the HumanEval benchmark~\cite{chen2021evaluating}. All 164 ground-truth specifications were manually written and cross-reviewed by researchers with multiple years of theorem-proving experience, requiring several person-weeks of expert effort; to the best of our knowledge, this is the first complete formal specification suite for HumanEval in Rocq. We use this suite as a reference baseline and evaluate a diverse set of state-of-the-art LLMs under a unified evaluation pipeline. Our evaluation clearly differentiates the specification generation capabilities of different models, revealing substantial performance gaps: \textsc{Coins} scores range from 28.05\% for Gemini~3 Pro Preview to 1.22\% for DeepSeek-V3.1. While frontier models can approach human-level performance on some tasks, generating tight formal specifications remains a significant challenge.



% Our results reveal substantial gaps between frontier and earlier-generation models and show that, while recent LLMs can approach human performance on some tasks, generating tight formal specifications remains a significant challenge.

Our main contributions are as follows:
\begin{itemize}
  \item We identify a fundamental limitation of existing program verification pipelines: while effective at checking whether implementations satisfy specifications, they provide limited ability to assess specification quality itself.
  \item We propose \textsc{Coins}, a principled evaluation framework for specification quality based on formally provable behavior in Rocq.
  \item We construct a curated dataset of human-written Rocq specifications for all 164 HumanEval problems and perform a comprehensive empirical evaluation of state-of-the-art LLMs.
  % \item Through extensive analysis, we provide new insights into the relationship between formal reasoning capability and specification quality.
\end{itemize}

